由整数除法算法与同余类代表的唯一性得到。平移群作用的轨道等于同余类，$\mathcal F_b$ 为其横截面，故投影 $\Pi_b$ 给出唯一 $r$，并据此确定 $q$。证毕。
\end{proof}

\begin{corollary}[中心余数横截面]\label{cor:pom-symmetric-remainder}
若改用对称基本域
\[
\mathcal F_b^{\mathrm{sym}}:=\left\{-\left\lfloor\frac{b}{2}\right\rfloor,\dots,\left\lceil\frac{b}{2}\right\rceil-1\right\},
\]
则仍有唯一分解 $a=qb+r$ 且 $r\in\mathcal F_b^{\mathrm{sym}}$。该选择等价于改变横截面而非改变本体演化，可将部分“区间校正”合并到单次序投影中。
\leanverified{symmetricQuoRem}
\leanverified{symmetricQuoRem\_eq}
\leanverified{symmetricQuoRem\_r\_symmetric}
\leanverified{paper\_pom\_symmetric\_remainder\_exists\_unique}
\end{corollary}

\begin{remark}[与 Fold 的群胚类比]
定理 \ref{thm:pom-division-section} 与定理 \ref{thm:zeckendorf-transversal} 在结构上同型：二者均为“轨道 = 值纤维/同余类，横截面 = 规范代表”的投影机制；商 $q$ 与折叠残差的群胚坐标具有同一语义地位。
\end{remark}

\subsection{结论簇 III：包含商余的扩展最小原语集}\label{subsec:pom-primitive-extended}
\begin{theorem}[扩展最小原语集]\label{thm:pom-extended-primitives}
\leanverified{paper\_pom\_extended\_primitives}
在配置层，$\oplus,\otimes,\mathrm{div},\mathrm{quo},\mathrm{rem}$ 可由以下原语生成：
\begin{enumerate}
  \item 逐点叠加（自由交换幺半群结构）；
  \item Fibonacci 群胚生成元（值保持局部可逆重写）；
  \item 值纤维横截面投影 $\Fold_\infty$（规范投影）；
  \item 平移群胚横截面投影 $\Pi_b$（序/区间投影）。
\end{enumerate}
在过程层，等价表述为：仅以复合 $\circ$ 作为底层原语，再配合值投影 $\Fold_\infty$ 与序投影 $\Pi_b$ 把输出落回外表；原子投影 $P_{\mathrm{prim}}$ 作为谱读出接口独立施加。
\end{theorem}
\begin{proof}
加法与乘法由定理 \ref{thm:composition-two-layer} 与 \ref{thm:arith-composition} 的过程层读出实现，并由 $\Fold_\infty$ 落回稳定地址。商余由定理 \ref{thm:pom-division-section} 作为平移群胚横截面投影给出。其余运算由上述原语的组合闭包得到。证毕。
\end{proof}

\begin{theorem}[双横截面正规形：值纤维 $\times$ 同余类]\label{thm:pom-double-transversal-normal-form}
对任意配置对象 $a\in\mathcal D$ 与任意整数 $b>0$，存在唯一三元组
\[
\boxed{(x,q,r)\in X_\infty\times\mathbb{Z}\times\mathcal F_b}
\]
使得
\[
x=\Fold_\infty(a),\qquad
r=\Pi_b(\mathrm{Val}(a)),\qquad
q=\frac{\mathrm{Val}(a)-r}{b},
\]
并且 \(\mathrm{Val}(a)=\mathrm{Val}(x)=qb+r\)。
\end{theorem}
\begin{proof}
由定理 \ref{thm:zeckendorf-transversal}（或等价的定义 \(\Fold_\infty(a)=Z(\mathrm{Val}(a))\)）得 \(x\) 唯一且满足 \(\mathrm{Val}(x)=\mathrm{Val}(a)\)。
由定理 \ref{thm:pom-division-section} 得 \((q,r)\) 唯一且满足 \(r=\Pi_b(\mathrm{Val}(a))\) 与 \(\mathrm{Val}(a)=qb+r\)。
合并两部分即得所示三元组的存在唯一性。证毕。
\leanpartial{paper\_pom\_double\_transversal\_normal\_form}{仅 $(q,r)$ 存在性 $\exists$；uniqueness 与 Fold 三元组 $x = \Fold_\infty(a)$ 待后轮}
\leanverified{Config}
\leanverified{val}
\leanverified{symRem}
\leanverified{symQuo}
\leanverified{symQuoRem\_spec}
\leanverified{symRem\_le}
\leanverified{symRem\_ge}
\leanverified{paper\_pom\_double\_transversal\_normal\_form\_unique\_standard}
\end{proof}

\begin{corollary}[投影调用的常数预算：末端正规化]\label{cor:pom-projection-budget}
设 \(t\) 为任意由配置层运算 \(\oplus,\otimes,\mathrm{div},\mathrm{quo},\mathrm{rem}\) 复合得到的表达式（深度任意），并假设其\emph{外部可见输出}仅为一次稳定地址读出（落在 \(X_\infty\)）以及可选的一次商余读出（落在 \(\mathbb Z\times\mathcal F_b\)）。
则存在一个等价的过程层实现，使得整条求值链中
\begin{itemize}
  \item 值横截面投影 \(\Fold_\infty\) 至多调用一次（仅在末端把输出落回 \(X_\infty\)）；
  \item 序横截面投影 \(\Pi_b\) 至多调用一次（仅当需要输出唯一余数代表时，于末端对 \(\mathrm{Val}\) 应用）。
\end{itemize}
并且：若输出类型要求唯一稳定地址（或要求唯一余数代表），则对应末端投影在该意义下不可省略。
\leanverified{paper\_pom\_projection\_budget}
\end{corollary}
\begin{proof}
由定理 \ref{thm:pom-extended-primitives}，过程层求值可改写为群胚内的可逆演化与末端横截面投影的复合。
而 \(\Fold_\infty\) 是 \(\mathrm{Val}\) 的规范截面：\(\mathrm{Val}(\Fold_\infty(a))=\mathrm{Val}(a)\)，且 \(\Fold_\infty(a)\) 仅依赖于 \(\mathrm{Val}(a)\)。
因此任何中间步骤对同一值纤维反复“先落回再继续演化”的正规化，都不会改变末端按最终值落回 \(X_\infty\) 的输出，可统一推迟为末端一次。
同理，\(\Pi_b\) 仅依赖于整数值 \(\mathrm{Val}(\cdot)\)，且 \(\mathrm{Val}\) 可在末端一次性读出，故若只需一次商余读出，则 \(\Pi_b\) 亦可推迟到末端一次。
最后，“不可省略性”来自轨道代表的非唯一性：不施加对应横截面投影，则输出仍停留在轨道上而非唯一点。证毕。
\end{proof}

\subsection{逆极限自然性与纤维不变量判别}\label{subsec:pom-new-results}
\subsubsection{跨分辨率一致性的 Gauge 原理：fold-aware restriction}\label{subsubsec:pom-gauge}
有限窗口微态空间 $\Omega_m=\{0,1\}^m$ 与稳定类型空间 $X_m$ 的角色并不对称：$X_m$ 在截断映射 $\pi_{m+1\to m}$ 下形成逆系统，而 $\Omega_m$ 在朴素截断下并不携带与折叠兼容的逆系统结构。该差异正是“跨分辨率一致性必须先规范化”的结构来源。

\begin{definition}[朴素截断与折叠感知限制]\label{def:pom-fold-aware-restriction}
设 $\tau_{m+1\to m}:\Omega_{m+1}\to\Omega_m$ 为朴素截断（丢弃最末一位；等价地取长度 $m$ 前缀）。定义折叠感知限制（fold-aware restriction）
$$
r^{\mathrm{e}}_{m+1\to m}:\Omega_{m+1}\to X_m,
\qquad
r^{\mathrm{e}}_{m+1\to m}:=\pi_{m+1\to m}\circ\Fold_{m+1}.
$$
\end{definition}

\begin{proposition}[朴素截断与折叠一般不交换]\label{prop:pom-truncation-not-commute}
\leanverified{paper\_pom\_truncation\_not\_commute}
一般地，图表
$$
\Omega_{m+1}\xrightarrow{\ \tau_{m+1\to m}\ }\Omega_m\xrightarrow{\ \Fold_m\ }X_m
$$
与
$$
\Omega_{m+1}\xrightarrow{\ \Fold_{m+1}\ }X_{m+1}\xrightarrow{\ \pi_{m+1\to m}\ }X_m
$$
并不一致，即存在 $\omega\in\Omega_{m+1}$ 使得
$$
\Fold_m(\tau_{m+1\to m}(\omega))\neq \pi_{m+1\to m}(\Fold_{m+1}(\omega)).
$$
\end{proposition}
\begin{proof}
折叠的局部重写（例如 $011\mapsto 100$）会产生跨位传播：被丢弃的末位可能决定某次局部重写是否触发，从而改变长度 $m$ 前缀在规范化后的形态。因此“先截断再折叠”与“先折叠再截断”在一般情形下不一致。证毕。
\end{proof}

\begin{proposition}[稳定逆系统的自然性：先折叠再限制]\label{prop:pom-gauge-natural}
对任意 $m\ge 1$，折叠感知限制满足
$$
r^{\mathrm{e}}_{m+1\to m}
=
\pi_{m+1\to m}\circ \Fold_{m+1},
$$
从而在稳定层形成严格可交换的跨分辨率图表
$$
\Omega_{m+1}\xrightarrow{\ \Fold_{m+1}\ }X_{m+1}\xrightarrow{\ \pi_{m+1\to m}\ }X_m
\qquad\text{与}\qquad
\Omega_{m+1}\xrightarrow{\ r^{\mathrm{e}}_{m+1\to m}\ }X_m
$$
一致；并且对任意 $m_2\ge m_1$ 有兼容性
$$
\text{令 }r^{\mathrm{e}}_{m_2\to m_1}:=\pi_{m_2\to m_1}\circ\Fold_{m_2},
\ \text{则对任意 }m_3\ge m_2\ge m_1\text{ 有 }r^{\mathrm{e}}_{m_3\to m_1}=\pi_{m_2\to m_1}\circ r^{\mathrm{e}}_{m_3\to m_2},
$$
因而“同一对象的跨分辨率坐标”只能在 $X_\infty\simeq\varprojlim X_m$ 的逆系统口径下被谈论。
\leanverified{paper\_pom\_gauge\_natural}
\end{proposition}
\begin{proof}
第一式为定义 \ref{def:pom-fold-aware-restriction} 的直接展开。兼容性由稳定层的逆系统映射满足的泛性质给出（定理 \ref{thm:pom-inverse-limit}），而折叠后的坐标才是逆系统对象。证毕。
\end{proof}

\begin{definition}[Gauge 曲率（2-cocycle obstruction）]\label{def:pom-gauge-curvature}
对给定的分辨率降阶 \((m{+}1\to m)\) 与顶层微态 \(\omega\in\Omega_{m+1}\)，定义曲率指示量
\[
\mathcal{K}_{m+1\to m}(\omega)
:=
\mathbf 1\!\left[\,
\Fold_m(\tau_{m+1\to m}(\omega))
\neq
\pi_{m+1\to m}(\Fold_{m+1}(\omega))
\,\right].
\]
它刻画“先截断再折叠”与“先折叠再限制”在该微态上的不交换障碍。
\end{definition}

\begin{proposition}[伪不变量的可测误差条：由曲率概率控制]\label{prop:pom-gauge-curvature-shadow-bound}
令 \(\mu\) 为 \(\Omega_{m+1}\) 上任意分布，令 \(f:X_m\to\mathbb{R}\) 为任意有界可观测量。
则有可审计的不等式
\[
\boxed{
\begin{aligned}
&\left|\mathbb E_{\omega\sim\mu}\,f(\Fold_m(\tau_{m+1\to m}(\omega)))
-\mathbb E_{\omega\sim\mu}\,f(\pi_{m+1\to m}(\Fold_{m+1}(\omega)))\right|\\
&\qquad \le 2\|f\|_\infty\cdot \mathbb P_{\omega\sim\mu}\!\left[\mathcal{K}_{m+1\to m}(\omega)=1\right].
\end{aligned}}
\]
特别地，在采用 fold-aware restriction 的 Gauge 口径下（命题 \ref{prop:pom-gauge-natural} 的严格交换图表），上述误差条塌缩为 \(0\)。
\end{proposition}
\begin{proof}
对每个 \(\omega\)，记
\(
x=\Fold_m(\tau_{m+1\to m}(\omega))
\)
与
\(
y=\pi_{m+1\to m}(\Fold_{m+1}(\omega))
\)。
则逐点有
\[
|f(x)-f(y)|
\le
2\|f\|_\infty\,\mathbf 1[x\neq y]
=
2\|f\|_\infty\,\mathcal{K}_{m+1\to m}(\omega).
\]
取 \(\omega\sim\mu\) 的期望即得所示不等式；最后一段由命题 \ref{prop:pom-gauge-natural} 直接推出。证毕。
\leanpartial{paper\_pom\_gauge\_curvature\_shadow\_bound}{仅形式化均匀 Finset 分布版本的耦合期望不等式 $|E[f(X)] - E[f(Y)]| \le 2M \cdot P[X \ne Y]$；任意 $\mu$ 与 $\mathrm{Fold}_m / \mathcal{K}$ 应用留后轮}
\leanverified{expectation}
\leanverified{probability}
\leanverified{neIndicator}
\leanverified{pointwise\_diff\_le}
\leanverified{abs\_expectation\_sub\_le}
\leanverified{probability\_eq\_expectation\_neIndicator}
\leanverified{coupling\_expectation\_bound}
\end{proof}

\begin{remark}[可复现审计：曲率概率的枚举/抽样]\label{rem:pom-gauge-curvature-audit}
右端的曲率概率
\(\mathbb P[\mathcal{K}_{m+1\to m}=1]\)
可直接由实现同时计算两条路径的输出并统计失配频率得到。
脚本 \path{scripts/exp_fold_truncation_curvature.py} 在均匀微态基线下对 \(m\le 20\) 做精确枚举并生成表格片段；对一般数据源亦可按同一定义以抽样估计。
例如，正文中用来定义 prime shadow 的计数函数、或任意 cylinder 指示函数 \(f\)，都是稳定类型的有界可观测量，因此不等式 \ref{prop:pom-gauge-curvature-shadow-bound} 给出一条可审计的“伪影误差条”。
\end{remark}

% Auto-generated by scripts/exp_fold_truncation_curvature.py
\begin{table}[H]
\centering
\small
\setlength{\tabcolsep}{8pt}
\renewcommand{\arraystretch}{1.12}
\begin{tabular}{rrrrrr}
\toprule
$m$ & $|\Omega_{m+1}|$ & $\mathbb P[\mathcal K_{m+1\to m}=1]$ & $\mathbb E[H_m]$ & $\mathbb E[H_m]/m$ & $\max H_m$ \\
\midrule
1 & 4 & 0.25 & 0.25 & 0.25 & 1 \\
2 & 8 & 0.25 & 0.25 & 0.125 & 1 \\
3 & 16 & 0.3125 & 0.375 & 0.125 & 2 \\
4 & 32 & 0.3125 & 0.40625 & 0.101562 & 3 \\
5 & 64 & 0.328125 & 0.453125 & 0.090625 & 4 \\
6 & 128 & 0.328125 & 0.46875 & 0.078125 & 5 \\
7 & 256 & 0.332031 & 0.484375 & 0.0691964 & 6 \\
8 & 512 & 0.332031 & 0.490234 & 0.0612793 & 7 \\
9 & 1024 & 0.333008 & 0.495117 & 0.055013 & 8 \\
10 & 2048 & 0.333008 & 0.49707 & 0.049707 & 9 \\
11 & 4096 & 0.333252 & 0.498535 & 0.0453214 & 10 \\
12 & 8192 & 0.333252 & 0.499146 & 0.0415955 & 11 \\
13 & 16384 & 0.333313 & 0.499573 & 0.0384287 & 12 \\
14 & 32768 & 0.333313 & 0.499756 & 0.0356968 & 13 \\
15 & 65536 & 0.333328 & 0.499878 & 0.0333252 & 14 \\
16 & 131072 & 0.333328 & 0.499931 & 0.0312457 & 15 \\
17 & 262144 & 0.333332 & 0.499966 & 0.0294097 & 16 \\
18 & 524288 & 0.333332 & 0.499981 & 0.0277767 & 17 \\
19 & 1048576 & 0.333333 & 0.49999 & 0.0263153 & 18 \\
20 & 2097152 & 0.333333 & 0.499995 & 0.0249997 & 19 \\
\bottomrule
\end{tabular}
\caption{Exact enumeration under the uniform baseline $\omega\sim\mathrm{Unif}(\Omega_{m+1})$ of the truncation--fold curvature indicator $\mathcal K_{m+1\to m}(\omega)=\mathbf 1[H_m(\omega)>0]$, where $H_m(\omega)$ is the Hamming distance between the two stable outputs $\Fold_m(\tau_{m+1\to m}(\omega))$ and $\pi_{m+1\to m}(\Fold_{m+1}(\omega))$.}
\label{tab:fold_truncation_curvature_stats}
\end{table}


\begin{lemma}[多步曲率由单步曲率事件支配]\label{lem:pom-multistep-curvature-dominated}
对任意整数 $n>m\ge 1$，令
\[
\tau_{n\to m}:=\tau_{m+1\to m}\circ\cdots\circ\tau_{n\to n-1},
\qquad
\pi_{n\to m}:=\pi_{m+1\to m}\circ\cdots\circ\pi_{n\to n-1},
\]
并定义多步不交换指示量
\[
\mathcal K_{n\to m}(\omega):=\mathbf 1\!\Big[\Fold_m(\tau_{n\to m}(\omega))\neq \pi_{n\to m}(\Fold_n(\omega))\Big]\qquad(\omega\in\Omega_n).
\]
则对所有 $\omega\in\Omega_n$ 成立
\[
\mathcal K_{n\to m}(\omega)\ \le\ \sum_{k=m}^{n-1}\mathcal K_{k+1\to k}\!\bigl(\tau_{n\to k+1}(\omega)\bigr).
\]
\end{lemma}
\begin{proof}
若右端为 $0$，则对每个 $k\in\{m,\dots,n-1\}$ 与 $\omega_{k+1}:=\tau_{n\to k+1}(\omega)\in\Omega_{k+1}$，由定义 $\mathcal K_{k+1\to k}(\omega_{k+1})=0$ 得
\[
\Fold_k(\tau_{k+1\to k}(\omega_{k+1}))=\pi_{k+1\to k}(\Fold_{k+1}(\omega_{k+1})).
\]
注意 $\tau_{k+1\to k}(\omega_{k+1})=\tau_{n\to k}(\omega)$。从 $k=m$ 起迭代并在每一步对等式两侧左复合相应的投影，利用逆系统的复合律 $\pi_{n\to m}=\pi_{k+1\to m}\circ\pi_{n\to k+1}$，可得
\[
\Fold_m(\tau_{n\to m}(\omega))=\pi_{n\to m}(\Fold_n(\omega)),
\]
即 $\mathcal K_{n\to m}(\omega)=0$。取逆否命题即得所示上界。证毕。
\end{proof}

\begin{theorem}[曲率事件可和性推出几乎处处的最终交换]\label{thm:pom-curvature-summable-eventual-commutation}
令 $\Omega_\infty:=\{0,1\}^{\NN}$，并以 $\tau_{\infty\to m}:\Omega_\infty\to\Omega_m$ 表示取前 $m$ 位前缀。
给定 $\Omega_\infty$ 上任意概率测度 $\PP$，定义事件
\[
A_m:=\Bigl\{\omega\in\Omega_\infty:\ \mathcal K_{m+1\to m}(\tau_{\infty\to m+1}(\omega))=1\Bigr\}.
\]
若满足可和性条件 $\sum_{m\ge 1}\PP(A_m)<\infty$，则存在 $\Omega_\infty^\star\subset\Omega_\infty$ 使 $\PP(\Omega_\infty^\star)=1$，且对每个 $\omega\in\Omega_\infty^\star$ 存在 $m_0(\omega)$ 使对所有 $m\ge m_0(\omega)$ 都有
\[
\Fold_m(\tau_{\infty\to m}(\omega))
=\pi_{m+1\to m}\!\Bigl(\Fold_{m+1}(\tau_{\infty\to m+1}(\omega))\Bigr).
\]
从而对每个 $\omega\in\Omega_\infty^\star$，存在唯一的逆极限点 $x_\infty(\omega)\in X_\infty\simeq\varprojlim X_m$ 满足
\[
\pi_m(x_\infty(\omega))=\Fold_m(\tau_{\infty\to m}(\omega))\qquad(m\ge m_0(\omega)).
\]
\end{theorem}
\leanverified{paper\_pom\_curvature\_summable\_eventual\_commutation}
\begin{proof}
由第一 Borel--Cantelli 引理，$\sum_m\PP(A_m)<\infty$ 蕴含 $\PP(A_m\ \mathrm{i.o.})=0$。因此对几乎处处的 $\omega$，存在 $m_0(\omega)$ 使得对所有 $m\ge m_0(\omega)$ 都有 $\omega\notin A_m$，即 $\mathcal K_{m+1\to m}(\tau_{\infty\to m+1}(\omega))=0$，展开定义得到所示最终交换恒等式。

令 $x_m(\omega):=\Fold_m(\tau_{\infty\to m}(\omega))$。对 $m\ge m_0(\omega)$，上式等价于 $\pi_{m+1\to m}(x_{m+1}(\omega))=x_m(\omega)$，从而 $(x_m(\omega))_{m\ge m_0(\omega)}$ 给出逆系统的一条相容尾部。对 $m<m_0(\omega)$ 再令 $x_m(\omega):=\pi_{m_0(\omega)\to m}(x_{m_0(\omega)}(\omega))$，得到全体 $m\ge 1$ 的相容族，于是由逆极限的泛性质对应唯一的 $x_\infty(\omega)\in X_\infty$。证毕。
\end{proof}

\begin{proposition}[曲率概率给出的总变差最优上界]\label{prop:pom-curvature-tv-optimal}
固定 $m\ge 1$，令 $\mu\in\Delta(\Omega_{m+1})$ 为任意分布，并在同一微态 $\omega\sim\mu$ 下定义
\[
X:=\Fold_m(\tau_{m+1\to m}(\omega))\in X_m,
\qquad
Y:=\pi_{m+1\to m}(\Fold_{m+1}(\omega))\in X_m,
\]
以及 $\varepsilon:=\PP_\mu[X\neq Y]=\PP_\mu[\mathcal K_{m+1\to m}(\omega)=1]$。
则
\[
D_{\mathrm{TV}}\bigl(\mathcal L(X),\mathcal L(Y)\bigr)\ \le\ \varepsilon,
\]
且该常数在一般情形下不可再改进。
\leanverified{paper\_pom\_curvature\_tv\_optimal}
\end{proposition}
\begin{proof}
对任意集合 $B\subseteq X_m$，有
\[
\bigl|\PP[X\in B]-\PP[Y\in B]\bigr|
=\bigl|\EE(\mathbf 1_B(X)-\mathbf 1_B(Y))\bigr|
\le \EE\,|\mathbf 1_B(X)-\mathbf 1_B(Y)|
\le \PP[X\neq Y]=\varepsilon.
\]
取上确界得到总变差上界。最优性可由支持几乎不交且 $\PP[X\neq Y]=\varepsilon$ 的构造实现，从而达到 $D_{\mathrm{TV}}=\varepsilon$。证毕。
\end{proof}

\begin{proposition}[曲率门槛对信息注入的可审计上界]\label{prop:pom-curvature-conditional-entropy-bound}
在命题 \ref{prop:pom-curvature-tv-optimal} 的设定下，令 $E:=\mathbf 1[X\neq Y]$，并记二元熵函数
\[
h(\varepsilon):=-\varepsilon\log\varepsilon-(1-\varepsilon)\log(1-\varepsilon),
\]
其中对数底任取但须一致。
则条件熵满足
\[
H(Y\mid X)\ \le\ h(\varepsilon)+\varepsilon\log\bigl(|X_m|-1\bigr).
\]
若进一步令末位比特 $B:=\omega_{m+1}\in\{0,1\}$，则
\[
I(B;Y\mid X)\ \le\ h(\varepsilon)+\varepsilon\log\bigl(|X_m|-1\bigr).
\]
\end{proposition}
\begin{proof}
在事件 $\{E=0\}$ 上有 $Y=X$，故 $H(Y\mid X,E=0)=0$。在事件 $\{E=1\}$ 上，给定 $X$ 时随机变量 $Y$ 的取值落在 $X_m\setminus\{X\}$，从而 $H(Y\mid X,E=1)\le \log(|X_m|-1)$。于是
\[
H(Y\mid X)\le H(E\mid X)+\PP(E=1)\,H(Y\mid X,E=1)
\le h(\varepsilon)+\varepsilon\log(|X_m|-1).
\]
又 $I(B;Y\mid X)=H(Y\mid X)-H(Y\mid X,B)\le H(Y\mid X)$，得到第二式。证毕。
\end{proof}

\begin{theorem}[截断--折叠曲率与单隐藏位的同一性]\label{thm:pom-truncation-curvature-equals-hidden-bit}
固定 \(m\ge 1\)，令 \(\eta\in\Omega_{m+1}\)。
定义曲率向量与指示量
\[
\kappa_{m+1\to m}(\eta)
:=\Fold_m(\tau_{m+1\to m}(\eta))\oplus \pi_{m+1\to m}(\Fold_{m+1}(\eta))\in\{0,1\}^m,
\qquad
\mathcal K_{m+1\to m}(\eta):=\mathbf 1[\kappa_{m+1\to m}(\eta)\neq 0],
\]
并在尺度 \(m{+}1\) 上以引理 \ref{lem:pom-one-bit} 的口径定义隐藏位 \(b(\eta)\in\{0,1\}\) 使得
\[
N(\eta)=V_{m+1}(\Fold_{m+1}(\eta))+b(\eta)\,F_{m+3}.
\]
则恒有
\[
\boxed{\;\mathcal K_{m+1\to m}(\eta)=b(\eta).\;}
\]
更精确地，令 \(r:=N(\eta)\bmod F_{m+2}\in\ZZ/(F_{m+2}\ZZ)\)，并沿用定理 \ref{thm:fold-congruence-universal-property} 的规范截面记号 \(w(r)\in X_m\)。
则
\[
\pi_{m+1\to m}(\Fold_{m+1}(\eta))=
\begin{cases}
w(r),& b(\eta)=0,\\
w(r-F_{m+1}),& b(\eta)=1,
\end{cases}
\qquad\text{从而}\qquad
\kappa_{m+1\to m}(\eta)=
\begin{cases}
0,& b(\eta)=0,\\
w(r)\oplus w(r-F_{m+1}),& b(\eta)=1.
\end{cases}
\]
\leanverified{paper\_pom\_truncation\_curvature\_equals\_hidden\_bit}
\end{theorem}
\begin{proof}
由引理 \ref{lem:pom-fold-congruence}（或定理 \ref{thm:fold-congruence-universal-property} 的等价表述），稳定输出 \(\Fold_m(\cdot)\) 仅依赖于 \(N(\cdot)\bmod F_{m+2}\)，并可写成 \(\Fold_m(\omega)=w(N(\omega)\bmod F_{m+2})\)。
朴素截断仅丢弃末位权重 \(\eta_{m+1}F_{m+2}\)，因此
\(
N(\tau_{m+1\to m}\eta)\equiv N(\eta)\pmod{F_{m+2}}
\)，从而
\[
\Fold_m(\tau_{m+1\to m}(\eta))=w(r).
\]
另一方面，对 \(\Fold_{m+1}\) 的隐藏位分解给出
\(
N(\eta)\equiv V_{m+1}(\Fold_{m+1}(\eta))\pmod{F_{m+3}}
\)，而把 \(\Fold_{m+1}(\eta)\) 投影到 \(X_m\) 并取 \(V_m\) 等价于在模 \(F_{m+2}\) 意义下忽略 \(F_{m+2}\) 位；因此
\[
V_m\!\left(\pi_{m+1\to m}(\Fold_{m+1}(\eta))\right)\equiv N(\eta)-b(\eta)F_{m+3}\pmod{F_{m+2}}.
\]
由 Fibonacci 递推 \(F_{m+3}=F_{m+2}+F_{m+1}\) 得 \(F_{m+3}\equiv F_{m+1}\ (\mathrm{mod}\ F_{m+2})\)，从而
\[
V_m\!\left(\pi_{m+1\to m}(\Fold_{m+1}(\eta))\right)\equiv r-b(\eta)F_{m+1}\pmod{F_{m+2}}.
\]
当 \(b(\eta)=0\) 时，上式与 \(V_m(w(r))=r\) 合并推出 \(\pi_{m+1\to m}(\Fold_{m+1}(\eta))=w(r)\)，故 \(\kappa_{m+1\to m}(\eta)=0\)，从而 \(\mathcal K_{m+1\to m}(\eta)=0\)。
当 \(b(\eta)=1\) 时，注意 \(F_{m+1}\not\equiv 0\ (\mathrm{mod}\ F_{m+2})\)，故余数类从 \(r\) 变为 \(r-F_{m+1}\)，推出
\(\pi_{m+1\to m}(\Fold_{m+1}(\eta))=w(r-F_{m+1})\neq w(r)\)，从而 \(\kappa_{m+1\to m}(\eta)=w(r)\oplus w(r-F_{m+1})\neq 0\)，即 \(\mathcal K_{m+1\to m}(\eta)=1\)。
两种情形合并得到 \(\mathcal K_{m+1\to m}(\eta)=b(\eta)\) 与所述精确表达式。证毕。
\end{proof}

\begin{proposition}[独立曲率尾部下的朴素坐标几乎处处失效]\label{prop:pom-curvature-independent-divergent}
在定理 \ref{thm:pom-curvature-summable-eventual-commutation} 的设定下仍记 $A_m=\{\mathcal K_{m+1\to m}(\tau_{\infty\to m+1}(\omega))=1\}$。
若进一步假设 $(A_m)_{m\ge 1}$ 相互独立且 $\sum_{m\ge 1}\PP(A_m)=\infty$，则 $\PP(A_m\ \mathrm{i.o.})=1$。
因此对 $\PP$-几乎处处的 $\omega$，序列 $x_m(\omega):=\Fold_m(\tau_{\infty\to m}(\omega))$ 违反相容条件 $\pi_{m+1\to m}(x_{m+1}(\omega))=x_m(\omega)$ 无穷多次，从而以朴素截断 $\tau_{\infty\to m}$ 作为跨分辨率坐标口径几乎处处不能产生逆极限坐标。
\end{proposition}
\begin{proof}
由第二 Borel--Cantelli 引理，在独立性与 $\sum_m\PP(A_m)=\infty$ 下有 $\PP(A_m\ \mathrm{i.o.})=1$。而 $A_m$ 的发生等价于 $\pi_{m+1\to m}(x_{m+1}(\omega))\neq x_m(\omega)$，故相容性失败无穷多次。证毕。
\end{proof}

\begin{corollary}[均匀微态基线下曲率概率的闭式与 \texorpdfstring{$1/3$}{1/3} 极限]\label{cor:pom-curvature-uniform-limit-closed}
令 \(\PP\) 为 \(\Omega_\infty\) 上的 i.i.d.\ Bernoulli\((\tfrac12)\) 乘积测度，并令
\[
\varepsilon_m:=\PP\!\left[\mathcal K_{m+1\to m}(\tau_{\infty\to m+1}(\omega))=1\right].
\]
则对所有 \(m\ge 1\) 有严格恒等式
\[
\boxed{\;
\varepsilon_m
=\PP[b(\eta)=1]_{\eta\sim\mathrm{Unif}(\Omega_{m+1})}
=\frac13-\frac{\varepsilon_{m+1}^\star}{3\cdot 2^{m+1}}\;},
\]
其中 \(\varepsilon_{m+1}^\star=1\)（\(m{+}1\) 偶）或 \(\varepsilon_{m+1}^\star=2\)（\(m{+}1\) 奇）。
因此
\[
\lim_{m\to\infty}\varepsilon_m=\frac13,\qquad
\varepsilon_m-\frac13=\Theta(2^{-m}).
\]
\end{corollary}
\begin{proof}
由定理 \ref{thm:pom-truncation-curvature-equals-hidden-bit} 有 \(\mathcal K_{m+1\to m}(\eta)=b(\eta)\)（\(\eta\in\Omega_{m+1}\)），故 \(\varepsilon_m=\PP[b(\eta)=1]\)（均匀微态）。
再由结论 \ref{con:pom-hidden-bit-mass-conservation-truncation} 的闭式计数（应用于分辨率 \(m{+}1\)）得到所示概率表达式与极限。证毕。
\end{proof}

\begin{proposition}[分辨率塔上的相容锥判别：曲率规约的必要性]\label{prop:pom-identity-cone-requires-curvature-reduction}
对任意 \(n>m\ge 1\) 与 \(\omega\in\Omega_n\)，定义缺陷向量
\[
D_{n\to m}(\omega):=\Fold_m(\tau_{n\to m}(\omega))\oplus \pi_{n\to m}(\Fold_n(\omega))\in\{0,1\}^m,
\qquad
\mathcal K_{n\to m}(\omega)=\mathbf 1[D_{n\to m}(\omega)\neq 0].
\]
令 \(\omega\in\Omega_\infty\)，并设 \(x_m(\omega):=\Fold_m(\tau_{\infty\to m}(\omega))\)。
则族 \((x_m(\omega))_{m\ge 1}\) 在逆系统中相容（即存在 \(x_\infty(\omega)\in\varprojlim X_m\) 使 \(\pi_m(x_\infty(\omega))=x_m(\omega)\) 对所有 \(m\) 成立）当且仅当对所有 \(m\ge 1\) 都有
\[
\mathcal K_{m+1\to m}(\tau_{\infty\to m+1}(\omega))=0.
\]
特别地，在均匀微态基线下由推论 \ref{cor:pom-curvature-uniform-limit-closed} 得 \(\PP[\mathcal K_{m+1\to m}=1]\to\tfrac13\)，因此朴素截断口径下的相容性失配并非可忽略的稀薄效应，而是由模结构不嵌套所强迫的稳定残差；相应的严格相容锥需以 fold-aware restriction 的交换口径（定义 \ref{def:pom-fold-aware-restriction}）给出。
\end{proposition}
\begin{proof}
相容性条件等价于对所有 \(m\) 有 \(\pi_{m+1\to m}(x_{m+1}(\omega))=x_m(\omega)\)。
按 \(x_m(\omega)\) 的定义展开即得
\(
\pi_{m+1\to m}(\Fold_{m+1}(\tau_{\infty\to m+1}\omega))=\Fold_m(\tau_{\infty\to m}\omega)
\)，这恰等价于 \(\mathcal K_{m+1\to m}(\tau_{\infty\to m+1}\omega)=0\)。
最后一段由推论 \ref{cor:pom-curvature-uniform-limit-closed} 给出的非零极限概率与定义 \ref{def:pom-fold-aware-restriction} 的严格交换图表合并得到。证毕。
\end{proof}

\begin{lemma}[朴素截断口径下缺陷的余圈组合律]\label{lem:pom-truncation-defect-cocycle}
在命题 \ref{prop:pom-identity-cone-requires-curvature-reduction} 的记号下，把 \(X_j\subseteq\{0,1\}^j\) 视为二元向量空间 \(V_j:=(\ZZ/2)^j\) 的子集，并将 \(\pi_{j\to i}\) 延拓为前缀投影 \(V_j\to V_i\)。
则对任意 \(n\ge k\ge m\ge 1\) 与任意 \(\omega\in\Omega_n\) 有严格恒等式
\[
\boxed{\;
D_{n\to m}(\omega)
\ =\
D_{k\to m}\!\bigl(\tau_{n\to k}(\omega)\bigr)\ \oplus\ \pi_{k\to m}\!\bigl(D_{n\to k}(\omega)\bigr)\; }.
\]
并且 \(D_{m\to m}\equiv 0\)。
\end{lemma}
\begin{proof}
取 \(\omega\in\Omega_n\)。
由定义
\[
D_{n\to m}(\omega)=\Fold_m(\tau_{n\to m}(\omega))\oplus \pi_{n\to m}(\Fold_n(\omega)).
\]
在 \(\ZZ/2\) 上插入并消去同一项 \(\pi_{k\to m}(\Fold_k(\tau_{n\to k}(\omega)))\)，并用 \(\tau_{n\to m}=\tau_{k\to m}\circ\tau_{n\to k}\)、\(\pi_{n\to m}=\pi_{k\to m}\circ\pi_{n\to k}\) 得
\[
\begin{aligned}
D_{n\to m}(\omega)
&=
\Bigl(\Fold_m(\tau_{k\to m}(\tau_{n\to k}\omega))\oplus \pi_{k\to m}(\Fold_k(\tau_{n\to k}\omega))\Bigr)\\
&\qquad\oplus
\pi_{k\to m}\Bigl(\Fold_k(\tau_{n\to k}\omega)\oplus \pi_{n\to k}(\Fold_n(\omega))\Bigr).
\end{aligned}
\]
第一括号按定义等于 \(D_{k\to m}(\tau_{n\to k}(\omega))\)，第二括号等于 \(\pi_{k\to m}(D_{n\to k}(\omega))\)，合并即得。
最后 \(D_{m\to m}(\omega)=\Fold_m(\omega)\oplus \Fold_m(\omega)=0\)。证毕。
\end{proof}

\begin{proposition}[缺陷余圈的 gauge 扭曲与余边界型变换]\label{prop:pom-truncation-defect-gauge-coboundary}
沿用上述记号。
给定任意 \(0\)-上链族 \(a=(a_m)_{m\ge 1}\)，其中 \(a_m:\Omega_m\to\{0,1\}^m\)。
定义 gauge 扭曲后的折叠族
\[
\Fold_m^{(a)}(\omega):=\Fold_m(\omega)\oplus a_m(\omega).
\]
令 \(D_{n\to m}^{(a)}\) 为由 \(\Fold^{(a)}\) 诱导的缺陷向量：
\[
D^{(a)}_{n\to m}(\omega):=\Fold_m^{(a)}(\tau_{n\to m}(\omega))\oplus \pi_{n\to m}(\Fold_n^{(a)}(\omega)).
\]
则对一切 \(n\ge m\) 与 \(\omega\in\Omega_n\) 有
\[
\boxed{\;
D_{n\to m}^{(a)}(\omega)
\ =\
D_{n\to m}(\omega)\ \oplus\ a_m(\tau_{n\to m}(\omega))\ \oplus\ \pi_{n\to m}\!\bigl(a_n(\omega)\bigr)\; }.
\]
\end{proposition}
\begin{proof}
直接代入定义并用 \(\pi_{n\to m}(u\oplus v)=\pi_{n\to m}(u)\oplus\pi_{n\to m}(v)\)（在 \(V_n\) 上）得到：
\[
\begin{aligned}
D^{(a)}_{n\to m}(\omega)
&=\bigl(\Fold_m(\tau_{n\to m}\omega)\oplus a_m(\tau_{n\to m}\omega)\bigr)\ \oplus\ 
\pi_{n\to m}\bigl(\Fold_n(\omega)\oplus a_n(\omega)\bigr)\\
&=D_{n\to m}(\omega)\ \oplus\ a_m(\tau_{n\to m}\omega)\ \oplus\ \pi_{n\to m}(a_n(\omega)).
\end{aligned}
\]
证毕。
\end{proof}

\begin{conjecture}[Zeckendorf 规约口径下的缺陷类非平凡性]\label{conj:pom-zeckendorf-defect-class-nontrivial}
设 \((\Fold_m)_{m\ge 1}\) 来自某个满足终止性、合流性且以 Zeckendorf 正规形为唯一正常形的二元局部重写实现（命题 \ref{prop:xi-local-binary-normalization-fibonacci-rigidity}）。
则朴素截断口径下的缺陷族 \((D_{n\to m})\)（引理 \ref{lem:pom-truncation-defect-cocycle}）在上述 gauge 扭曲意义下不平凡：不存在 \(0\)-上链族 \(a=(a_m)\) 使得对所有 \(n\ge m\) 恒有 \(D^{(a)}_{n\to m}\equiv 0\)。
等价地，缺陷余圈的 gauge 等价类张成一个非零的 \(\ZZ_2\)-一维子空间，并且仅依赖于 Zeckendorf 正规形体系而与具体重写策略无关。
\end{conjecture}

\begin{remark}[曲率尾部门槛与逆极限粒子化判据]\label{rem:pom-curvature-particleization}
跨分辨率比较若以朴素前缀 $\tau$ 为坐标口径，则曲率事件控制了“先截断再折叠”与“先折叠再限制”的不一致频率，从而把伪不变量风险转化为可审计的尾部门槛。
在给定微态基线 $\PP$ 下，可将可讨论的稳定本体限定为满足曲率尾部可和（或最终为零）的逆极限子集：其不同分辨率坐标 $\pi_m(x_\infty)$ 才能被视为同一对象在稳定逆系统口径下的一致投影。
\end{remark}

\begin{proposition}[\texorpdfstring{$\mathbb Z_2$}{Z2}-障碍类口径下的“曲率--进位缺陷”同构原型]\label{prop:pom-curvature-carry-defect-z2-obstruction-principle}
定义 \ref{def:fold-local-curvature-defect} 的局部交换缺陷向量 \(\kappa_{m+1\to m}(\eta)\in\{0,1\}^m\) 与其指示 \(K_{m+1\to m}(\eta)\in\{0,1\}\)（亦见定理 \ref{thm:pom-truncation-curvature-equals-hidden-bit} 的 \(\mathcal K_{m+1\to m}\) 记号）给出一个以 \(\mathbb Z_2\) 为系数的障碍型不变量：其在尺度塔上的累积满足离散广义 Stokes 恒等式（定理 \ref{thm:fold-discrete-stokes-defect}），因此以“异或求和”方式实现可加性。
\par
另一方面，附录定理 \ref{thm:multinomial-vp-carry-signature} 将多项式系数的 \(p\)-进赋值识别为 \(p\)-进逐位加法的总进位次数；特别在 \(p=2\) 时，该总进位次数在 \(\mathbb Z_2\) 上同样具有可加性，并在折叠诱导的规范体积链上表现为尺度缺陷（定理 \ref{thm:fold-multiscale-carry-defect}）。
\par
因此在 \(\mathbb Z_2\)-上同调的抽象口径下，“折叠--截断的不交换曲率”与“二进进位缺陷”可视为同一类对象的两种实例：二者均由逐位进位/溢出机制生成，并作为取值于 \(\mathbb Z_2\) 的障碍类刻画跨尺度可交换性的失配。
\end{proposition}


\begin{definition}[分辨率反射子（跨分辨率纤维均匀化）]\label{def:pom-resolution-reflector}
固定顶层分辨率 $M$，并把微态空间取为 $\Omega_M$。
对任意 $1\le m\le M$，以折叠感知限制
\(
r^{\mathrm{e}}_{M\to m}=\pi_{M\to m}\circ \Fold_M:\Omega_M\to X_m
\)
为 coarse-graining 映射，定义推前算子
\[
P_{M\to m}:=(r^{\mathrm{e}}_{M\to m})_*:\Delta(\Omega_M)\to\Delta(X_m),
\]
并记其纤维基数为
\[
d_{M\to m}(x):=\bigl|(r^{\mathrm{e}}_{M\to m})^{-1}(x)\bigr|\qquad(x\in X_m).
\]
定义纤维均匀提升算子
\[
Q_{M\to m}:\Delta(X_m)\to\Delta(\Omega_M),
\qquad
(Q_{M\to m}\pi)(\omega):=\frac{\pi(r^{\mathrm{e}}_{M\to m}(\omega))}{d_{M\to m}(r^{\mathrm{e}}_{M\to m}(\omega))}.
\]
则 $P_{M\to m}\circ Q_{M\to m}=\mathrm{id}_{\Delta(X_m)}$。
对任意 $\mu\in\Delta(\Omega_M)$，定义其分辨率 $m$ 的规范化反射（纤维均匀化）为
\[
\mu^{(m)}:=Q_{M\to m}P_{M\to m}\mu\in\Delta(\Omega_M).
\]
当 $m=M$ 时，上述构造退化为小节 \ref{subsec:pom-information-ledger} 的 $P_M,Q_M$ 与 $\mu^{e}$（仅记号不同）。
\end{definition}

\begin{theorem}[KL-毕达哥拉斯：分辨率塔上的残差可加]\label{thm:pom-kl-pythagoras-tower}
\leanverified{paper\_pom\_kl\_pythagoras\_tower}
沿用定义 \ref{def:pom-resolution-reflector}，并以定理 \ref{thm:pom-kl-ledger} 的口径记 KL 散度。
则对任意 $M\ge m_2\ge m_1\ge 1$ 与任意 $\mu\in\Delta(\Omega_M)$，成立严格恒等式
\[
\boxed{\;
D_{\mathrm{KL}}(\mu\|\mu^{(m_1)})
=D_{\mathrm{KL}}(\mu\|\mu^{(m_2)})
+D_{\mathrm{KL}}(\mu^{(m_2)}\|\mu^{(m_1)})\; }.
\]
\end{theorem}
\begin{proof}
由 KL 定义，
\[
D_{\mathrm{KL}}(\mu\|\mu^{(m_1)})-D_{\mathrm{KL}}(\mu\|\mu^{(m_2)})
=\sum_{\omega\in\Omega_M}\mu(\omega)\log\frac{\mu^{(m_2)}(\omega)}{\mu^{(m_1)}(\omega)}.
\]
对固定的 $y\in X_{m_2}$，在纤维 $(r^{\mathrm{e}}_{M\to m_2})^{-1}(y)$ 上，上式中的比值
\(
\mu^{(m_2)}(\omega)/\mu^{(m_1)}(\omega)
\)
为常数（因为 $\mu^{(m)}$ 在每条 $r^{\mathrm{e}}_{M\to m}$-纤维上均匀）。
因此按 $y$ 分组可得
\[
\sum_{\omega\in\Omega_M}\mu(\omega)\log\frac{\mu^{(m_2)}(\omega)}{\mu^{(m_1)}(\omega)}
=\sum_{y\in X_{m_2}}(P_{M\to m_2}\mu)(y)\,
\log\frac{\mu^{(m_2)}(\omega_y)}{\mu^{(m_1)}(\omega_y)},
\]
其中 $\omega_y$ 为该纤维内任意代表点。
另一方面，由定义 \(\mu^{(m_2)}=Q_{M\to m_2}P_{M\to m_2}\mu\) 得
\(
P_{M\to m_2}\mu^{(m_2)}=P_{M\to m_2}\mu
\)，
从而上式右端可将系数 $(P_{M\to m_2}\mu)(y)$ 替换为 $(P_{M\to m_2}\mu^{(m_2)})(y)$。
再将分组表达“反展开”回 $\Omega_M$，即得
\[
\sum_{\omega\in\Omega_M}\mu(\omega)\log\frac{\mu^{(m_2)}(\omega)}{\mu^{(m_1)}(\omega)}
=\sum_{\omega\in\Omega_M}\mu^{(m_2)}(\omega)\log\frac{\mu^{(m_2)}(\omega)}{\mu^{(m_1)}(\omega)}
=D_{\mathrm{KL}}(\mu^{(m_2)}\|\mu^{(m_1)}).
\]
移项即得结论。证毕。
\end{proof}

\begin{corollary}[跨分辨率的信息守恒账本：望远镜分解]\label{cor:pom-kl-ledger-telescope}
\leanverified{paper\_pom\_kl\_ledger\_telescope}
在定理 \ref{thm:pom-kl-pythagoras-tower} 的设定下，定义单步外置量
\[
\Delta_m(\mu):=D_{\mathrm{KL}}(\mu^{(m)}\|\mu^{(m-1)})\qquad(2\le m\le M).
\]
则对任意 $M\ge 1$ 有严格望远镜分解
\[
\boxed{\;
D_{\mathrm{KL}}(\mu\|\mu^{(1)})
=\sum_{m=2}^{M}\Delta_m(\mu)\;+\;D_{\mathrm{KL}}(\mu\|\mu^{(M)})\; }.
\]
\end{corollary}
\begin{proof}
对 $m=2,3,\dots,M$ 依次将定理 \ref{thm:pom-kl-pythagoras-tower} 应用于三元组 $(m_1,m_2)=(m-1,m)$ 并求和，所有中间项望远镜抵消，得到所示恒等式。证毕。
\end{proof}

\begin{remark}[结构判据：伪不变量来自坐标不一致]\label{rem:pom-gauge-criterion}
若某一 ISA/微码/核在跨分辨率上使用朴素截断（$\tau$）来比较对象，而非经 $P_Z$（或等价地经 $r^{\mathrm{e}}$）落到稳定逆系统，则其“规律/不变量”可能只是坐标不一致产生的幻影：同一高分辨率微态在低分辨率下被错误归类，从而制造伪素数影子与伪收敛现象。因而跨分辨率可审计性必须以 fold-aware restriction 作为 Gauge 条件。
\end{remark}

\begin{theorem}[逆极限自然性]\label{thm:pom-inverse-limit}
\leanverified{paper\_pom\_inverse\_limit}
设对每个分辨率 $m$ 定义稳定层运算 $f_m:X_m\to X_m$，并满足对一切 $m_2\ge m_1$ 的截断投影
\[
\pi_{m_2\to m_1}\circ f_{m_2}=f_{m_1}\circ \pi_{m_2\to m_1}.
\]
则存在唯一的极限运算 $f_\infty:X_\infty\to X_\infty$ 使得对所有 $m$ 有 $\pi_m\circ f_\infty=f_m\circ \pi_m$。
\end{theorem}
\begin{proof}
由定理 \ref{thm:inverse-limit-golden}，$X_\infty\cong\varprojlim X_m$。相容性等式正是逆系统上的自然变换条件，因此由逆极限的泛性质得到唯一的 $f_\infty$。证毕。
\end{proof}

\begin{proposition}[纤维不变量判别]\label{prop:pom-fiber-invariance}
\leanverified{paper\_pom\_fiber\_invariance}
给定 $g:\Omega_m\to\mathbb{R}$，以下条件等价：
\begin{enumerate}
  \item 若 $\Fold_m(a)=\Fold_m(b)$ 则 $g(a)=g(b)$；
  \item 存在 $\bar g:X_m\to\mathbb{R}$ 使得 $g=\bar g\circ\Fold_m$；
  \item $g$（作为可观测量）属于折叠不变量子代数 $\mathcal{B}_m$（命题 \ref{prop:fold-invariant-subalgebra}）。
\end{enumerate}
\end{proposition}
\begin{proof}
$(1)\Rightarrow(2)$：定义 $\bar g(x)=g(a)$ 对任意 $a\in\Fold_m^{-1}(x)$，由 (1) 保证良定义。$(2)\Rightarrow(3)$：$g$ 在每条纤维上常值，故对应于 $\mathcal{B}_m$ 的坐标表示。$(3)\Rightarrow(1)$：$\mathcal{B}_m$ 的元素对纤维常值（命题 \ref{prop:fold-invariant-subalgebra}）。证毕。
\end{proof}

\begin{remark}[时间—空间对偶接口与素数影子]\label{rem:pom-time-space}
命题 \ref{prop:time-space-commuting-diagram} 给出的严格交换图把时间侧 $\Tr(T^n)$ 与空间侧 primitive 轨道谱对齐，从而为“时间—空间正交对偶”提供可审计接口。另一方面，有限核的 primitive 轨道谱与整数素数并不等同（见注 \ref{rem:prime-shadow-vs-primes}）：前者是核动力学的内禀素性，后者需要额外的素数索引算子族。素数寄存器纤维 $\mathcal{P}$ 正是对这一额外索引结构的本体化补充。
\end{remark}

\begin{theorem}[profinite 轴上的 Chebotarev 提升]\label{thm:pom-profinite-axis}
\leanverified{paper\_pom\_profinite\_axis}
设 $\{G_m,\pi_{m_2\to m_1}\}$ 为有限群的投影系统，$G_\infty:=\varprojlim_m G_m$。对每个 $m$，考虑对应的 $G_m$-扩张核，并记
$$
\lambda_m:=\rho(B_{\mathbf 1,m}),\qquad
\Lambda_m:=\max_{\rho\neq\mathbf 1}\rho(B_{\rho,m}).
$$
令 \(U\subset G_\infty\) 为 cylinder 集合，即存在某个 $m$ 与共轭不变子集 $A\subseteq G_m$ 使
\(
U=\pi_m^{-1}(A)
\)。
则对应的 primitive 轨道计数满足
$$
B_{n,U}
=\mu_{\mathrm{Haar}}(U)\,\frac{\lambda_m^n}{n}
 +O\!\left(\frac{\Lambda_m^n}{n}\right).
$$
其中 \(\mu_{\mathrm{Haar}}(U)=|A|/|G_m|\)。
若进一步满足核版 GRH 条件 $\Lambda_m\le \lambda_m^{1/2}$（对该层 $m$），则误差降为平方根级
$$
B_{n,U}
=\mu_{\mathrm{Haar}}(U)\,\frac{\lambda_m^n}{n}
 +O\!\left(\frac{\lambda_m^{n/2}}{n}\right).
$$
\end{theorem}
\begin{proof}
由于 $U$ 为 cylinder，存在 $m$ 与共轭不变子集 $A\subseteq G_m$ 使 $U=\pi_m^{-1}(A)$。令 $B_{n,c}$ 表示落在共轭类 $c\subseteq G_m$ 的 primitive 轨道计数，则
$$
B_{n,U}=\sum_{c\subseteq A} B_{n,c}.
$$
对每个 $c$，由有限层的 Chebotarev 型估计（定理 \ref{thm:kernel-chebotarev-exp}）有
$
B_{n,c}=\frac{|c|}{|G_m|}\frac{\lambda_m^n}{n}+O(\Lambda_m^n/n)
$。
将各类相加得到主项系数
$
\sum_{c\subseteq A}|c|/|G_m|=|A|/|G_m|=\mu_{\mathrm{Haar}}(U)
$，并因 $A$ 的共轭类数目有限，误差仍为 $O(\Lambda_m^n/n)$。若 $\Lambda_m\le \lambda_m^{1/2}$，则立得平方根级误差。证毕。
\end{proof}

\begin{corollary}[cylinder 级有效均匀化长度（相对误差的对数预算）]\label{cor:pom-profinite-cylinder-effective-n}
沿用定理 \ref{thm:pom-profinite-axis} 的记号。令 \(U=\pi_m^{-1}(A)\subseteq G_\infty\) 为一个层级为 \(m\) 的 cylinder，且假设 \(A\neq\varnothing\)。
则对长度 \(n\) 的 primitive 轨道计数，有相对误差界
\[
\frac{\bigl|B_{n,U}-\mu_{\mathrm{Haar}}(U)\,\lambda_m^n/n\bigr|}{\mu_{\mathrm{Haar}}(U)\,\lambda_m^n/n}
\;=\;
O\!\left(\frac{|G_m|}{|A|}\left(\frac{\Lambda_m}{\lambda_m}\right)^n\right)
\;\le\;
O\!\left(|G_m|\left(\frac{\Lambda_m}{\lambda_m}\right)^n\right).
\]
特别地，若存在 \(\varepsilon>0\) 使得该层满足统一谱隙形式
\[
\Lambda_m\ \le\ \lambda_m^{1/2-\varepsilon},
\]
则对任意目标相对误差阈值 \(\delta\in(0,1)\)，只要
\[
\boxed{\;
n\ \gtrsim\ \frac{\log(|G_m|/\delta)}{(1/2+\varepsilon)\log \lambda_m}\; }
\]
即可保证该 cylinder 上的 Chebotarev 主项达到相对误差 \(O(\delta)\) 的尺度。
\leanverified{paper\_pom\_profinite\_cylinder\_effective\_n}
\end{corollary}

\begin{proof}
由定理 \ref{thm:pom-profinite-axis}，
\(
|B_{n,U}-\mu_{\mathrm{Haar}}(U)\lambda_m^n/n|\le C\,\Lambda_m^n/n
\)
（其中常数 \(C\) 仅依赖于该层的有限核与 \(A\) 的共轭类个数）。
两边除以主项 \(\mu_{\mathrm{Haar}}(U)\lambda_m^n/n=(|A|/|G_m|)\lambda_m^n/n\) 得
\[
\frac{|B_{n,U}-\mu_{\mathrm{Haar}}(U)\lambda_m^n/n|}{\mu_{\mathrm{Haar}}(U)\lambda_m^n/n}
\le
\frac{C|G_m|}{|A|}\left(\frac{\Lambda_m}{\lambda_m}\right)^n.
\]
若进一步 \(\Lambda_m\le \lambda_m^{1/2-\varepsilon}\)，则 \((\Lambda_m/\lambda_m)^n\le \lambda_m^{-(1/2+\varepsilon)n}\)。
令右端 \(\le \delta\) 并解出 \(n\) 即得所示对数预算。证毕。
\end{proof}

\begin{corollary}[反解：给定时间 \(n\) 的可承载 cylinder/寄存器容量上界]\label{cor:pom-profinite-cylinder-capacity}
沿用推论 \ref{cor:pom-profinite-cylinder-effective-n} 的记号与统一谱隙假设
\(
\Lambda_m\le \lambda_m^{1/2-\varepsilon}
\)。
固定目标相对误差阈值 \(\delta\in(0,1)\) 与长度预算 \(n\)。
则要在层级 \(m\) 的任意非空 cylinder \(U=\pi_m^{-1}(A)\)（\(A\neq\varnothing\)）上把 Chebotarev 相对误差压到 \(O(\delta)\) 的尺度，一个（在 \(\lesssim\) 口径下）足够的容量条件是
\[
\boxed{\;
\log |G_m|\ \lesssim\ (1/2+\varepsilon)\,n\log\lambda_m\ -\ \log(1/\delta)\; }.
\]
换言之：在统一谱隙口径下，profinite 轴上“可承载的寄存器容量” \(\log|G_m|\) 与“时间步长” \(n\) 之间呈线性上界。
\leanverified{paper\_pom\_profinite\_cylinder\_capacity}
\end{corollary}
\begin{proof}
由推论 \ref{cor:pom-profinite-cylinder-effective-n} 的推导过程（见其证明中的显式界）
\(
\mathrm{Err}\lesssim |G_m|\,\lambda_m^{-(1/2+\varepsilon)n}
\)。
令右端 \(\le \delta\) 并取对数，即得到所示容量上界（把隐含常数吸收进 \(\lesssim\)）。
证毕。
\end{proof}

\begin{theorem}[同步扣除的 Chebotarev 容量律]\label{thm:pom-sync-subtracted-chebotarev-capacity}
\leanverified{paper\_pom\_sync\_subtracted\_chebotarev\_capacity}
固定层级 \(m\) 与两类置信参数 \(\delta_{\mathrm{NULL}},\delta_{\mathrm{Cheb}}\in(0,1)\)。
沿用同步核的右解析确定化表示与同步时间 \(T_m\)（定义 \ref{def:Ym-sync-time}），并按“地址先于概念值/未同步即无地址”的语义约定把 \(T_m\) 解释为统计型 \(\mathsf{NULL}\) 的消失时刻（备注 \ref{rem:Ym-stat-null-unsync}）。
\par
定义同步扣除项
\[
n_{\mathrm{sync}}(m,\delta_{\mathrm{NULL}})
:=
\begin{cases}
D_{\mathrm{sync}}(m), & \text{若 }Y_m^{\mathrm{amb}}=\varnothing,\\[6pt]
\left\lceil \dfrac{1}{\varepsilon_m}\log\dfrac{C_m}{\delta_{\mathrm{NULL}}}\right\rceil,
& \text{若 }Y_m^{\mathrm{amb}}\neq\varnothing,
\end{cases}
\]
其中 \(D_{\mathrm{sync}}(m)\) 由推论 \ref{cor:Ym-sync-delay}（或更一般由备注 \ref{rem:Ym-amb-empty} 的 DAG 最长路审计口径）给出；
而第二行的常数 \(C_m,\varepsilon_m\) 来自定理 \ref{thm:Ym-sync-tail} 的指数尾预算。
并令可用深度
\[
n_{\mathrm{eff}}:=\bigl(n-n_{\mathrm{sync}}(m,\delta_{\mathrm{NULL}})\bigr)_+,
\qquad (x)_+:=\max\{x,0\}.
\]
\par
若该层满足统一谱隙假设
\(
\Lambda_m\le \lambda_m^{1/2-\varepsilon_{\mathrm{gap}}}
\)
（对某个 \(\varepsilon_{\mathrm{gap}}>0\)），则任何企图在总扫描预算 \(n\) 内同时满足
\begin{enumerate}
  \item \textbf{读出定义性：}\(\nu_m(T_m>n_{\mathrm{sync}})\le \delta_{\mathrm{NULL}}\)（从而地址以概率 \(\ge 1-\delta_{\mathrm{NULL}}\) 被前缀唯一化），以及
  \item \textbf{读出统计清晰度：}对任意非空 cylinder \(U=\pi_m^{-1}(A)\) 把 Chebotarev 相对误差压到 \(\le O(\delta_{\mathrm{Cheb}})\)，
\end{enumerate}
都必须把 \(n_{\mathrm{eff}}\) 作为进入 Chebotarev 预算的“可用深度”。等价地，推论 \ref{cor:pom-profinite-cylinder-capacity} 的容量条件发生同步扣除修正：
\[
\boxed{\;
\log |G_m|
\ \lesssim\
(1/2+\varepsilon_{\mathrm{gap}})\,n_{\mathrm{eff}}\log\lambda_m
\ -\
\log(1/\delta_{\mathrm{Cheb}})
\;}
\]
（隐含常数口径与推论 \ref{cor:pom-profinite-cylinder-capacity} 一致）。
\par
\noindent\textbf{可审计性说明（作为结论的一部分）.}
\begin{enumerate}
  \item 同步扣除项 \(n_{\mathrm{sync}}\) 在有限确定化图上可直接审计：歧义子图含有向环当且仅当 \(Y_m^{\mathrm{amb}}\neq\varnothing\)，无环则 \(D_{\mathrm{sync}}\) 由最长路 DP 得到（备注 \ref{rem:Ym-amb-empty} 与表 \ref{tab:phi_m_sync_theory}）。
  \item 同步延迟属于不可空间化时间不变量：它与 residue/prime-register 等正交记录轴的扩展正交，不能通过单纯增加模轴容量消去；因此必须以显式时间扣除项进入任何容量预算（见备注 \ref{rem:pom-time-defect-dichotomy}）。
\end{enumerate}
\end{theorem}

\begin{corollary}[Zeckendorf 折叠下的尖锐线性扣除]\label{cor:pom-sync-subtracted-chebotarev-zeckendorf}
\leanverified{paper\_pom\_sync\_subtracted\_chebotarev\_zeckendorf}
在本文的 Zeckendorf 折叠模型中，当 \(m\ge 3\) 时有 \(Y_m^{\mathrm{amb}}=\varnothing\) 且 \(D_{\mathrm{sync}}(m)=m-1\)（推论 \ref{cor:Ym-sync-delay} 与备注 \ref{rem:Ym-amb-empty}）。
因此对任意 \(\delta_{\mathrm{NULL}},\delta_{\mathrm{Cheb}}\in(0,1)\)，有
\[
n_{\mathrm{eff}}=(n-m+1)_+,
\]
从而定理 \ref{thm:pom-sync-subtracted-chebotarev-capacity} 的容量条件退化为尖锐的线性扣除形式
\[
\log |G_m|
\ \lesssim\
(1/2+\varepsilon_{\mathrm{gap}})\,(n-m+1)_+\log\lambda_m
\ -\
\log(1/\delta_{\mathrm{Cheb}}).
\]
\end{corollary}

\begin{corollary}[弱可逆情形的两参数预算]\label{cor:pom-sync-subtracted-chebotarev-weak}
若 \(Y_m^{\mathrm{amb}}\neq\varnothing\)，则为使统计型 \(\mathsf{NULL}\) 在最大熵测度下以概率 \(\le \delta_{\mathrm{NULL}}\) 消失，必须先支付同步预算
\[
n_{\mathrm{sync}}(m,\delta_{\mathrm{NULL}})
=\left\lceil \dfrac{1}{\varepsilon_m}\log\dfrac{C_m}{\delta_{\mathrm{NULL}}}\right\rceil,
\]
之后才存在统一的 Chebotarev 统计预算可谈；并且时间侧容量条件以 \(n_{\mathrm{eff}}=(n-n_{\mathrm{sync}})_+\) 的扣除方式耦合两类指数率：\(\varepsilon_m\) 控制定义性清晰度的同步尾速率，而 \(\lambda_m\) 与 \(\varepsilon_{\mathrm{gap}}\) 控制统计清晰度的谱速率（定理 \ref{thm:pom-sync-subtracted-chebotarev-capacity}）。
\end{corollary}
\leanverified{paper\_pom\_sync\_subtracted\_chebotarev\_weak}

\subsubsection{\texorpdfstring{$L^2(\widehat{\ZZ})$}{L2(Zhat)} 上的寄存器提升算子：角色直和分解与“统一全部模数”的算子框架}\label{subsec:profinite-operator-l2}
定理 \ref{thm:pom-profinite-axis} 以“每个有限层 \(G_m\)”为单位给出 Chebotarev 计数与误差口径。
本小节给出一个更接近 Hilbert--P{\'o}lya 语义的压缩：把所有模数的扭曲通道统一为一个作用在 \(L^2(G_\infty)\) 上的提升算子，并在 Pontryagin 对偶侧把它严格对角化为一族 twist 矩阵。

\paragraph{专门化到素数寄存器轴：\texorpdfstring{$G_\infty=\widehat{\ZZ}$}{G=Zhat}}
取
\[
G_\infty=\widehat{\ZZ}=\varprojlim_m \ZZ/m\ZZ,
\]
它是紧阿贝尔群，其 Pontryagin 对偶为离散群
\[
\widehat{G_\infty}\cong \QQ/\ZZ,
\]
并存在一组标准正交角色基 \(\{\chi_\alpha\}_{\alpha\in\QQ/\ZZ}\subset L^2(\widehat{\ZZ})\)；见正文关于角色基与对偶的讨论，以及附录中同余轴的 DFT 口径。

\paragraph{提升算子}
设核的状态集合为有限集 \(V\)，并记带权有向边集合为 \(E\subset V\times V\)。
对每条边 \(e:i\to j\) 赋一个权重 \(w(e)\in\CC\)（例如由输出位势/碰撞势/温度化门给出），并赋一个“寄存器增量”标记
\(\gamma(e)\in \widehat{\ZZ}\)
（在有限层 \(\ZZ/m\ZZ\) 上它就是同余更新步；逆极限上由相容性唯一确定）。
定义 Hilbert 空间
\[
\mathcal{H}:=\ell^2(V)\otimes L^2(\widehat{\ZZ}),
\]
并定义提升算子 \(T:\mathcal{H}\to\mathcal{H}\) 为
\[
\boxed{\ (Tf)_j(x):=\sum_{e:i\to j} w(e)\,f_i\!\bigl(x-\gamma(e)\bigr)\qquad(x\in\widehat{\ZZ})\ }.
\]
这正是“有限状态 \(\times\) 紧群平移”型的标准构造：它把有限核的转移与寄存器平移耦合为一个单一算子。

\paragraph{角色分解：一切扭曲通道同时对角化}
由于 \(\widehat{\ZZ}\) 为紧阿贝尔群，Fourier 变换给出幺正同构
\[
L^2(\widehat{\ZZ})\ \cong\ \ell^2(\QQ/\ZZ),
\]
并把平移作用对角化为角色乘子。于是 \(T\) 在 Fourier 侧分解为一个\textbf{严格直和}（可视为以计数测度的直积分）：
\[
\boxed{\ T\ \simeq\ \bigoplus_{\alpha\in\QQ/\ZZ} M_{\chi_\alpha}\ },
\]
其中每个纤维 twist 矩阵 \(M_{\chi}\in\CC^{V\times V}\) 的条目为
\[
\boxed{\ (M_\chi)_{ij}:=\sum_{e:i\to j} w(e)\,\chi\!\bigl(\gamma(e)\bigr)\ }.
\]
因此“对所有模数/所有角色的扭曲”被统一为同一个算子 \(T\) 的角色分量谱。
在有限层 \(G_m=\ZZ/m\ZZ\) 上，上式退化为离散傅里叶变换（DFT）下的块对角化，正是定理 \ref{thm:pom-profinite-axis} 与有限层 Chebotarev 估计背后的线性代数机制。

\begin{remark}[Grand Unified Lift：把通道/边界塔并入提升算子纤维]\label{rem:pom-grand-unified-lift}
上文的 \(T\) 统一的是“同余/模数轴”上的全部扭曲通道。
若进一步希望把 \((\varphi,\pi,e)\) 三投影通道与多分辨率边界塔一起并入同一算子框架，
可取一个额外的紧群（或有限群）纤维 \(G_{\mathrm{chan}}\) 来承载这些离散标签，并定义扩展纤维
\[
G_\infty^{\mathrm{GUT}}
:=\widehat{\ZZ}\times G_{\mathrm{chan}},
\qquad
\mathcal{H}_{\mathrm{GUT}}:=\ell^2(V)\otimes L^2\!\left(G_\infty^{\mathrm{GUT}}\right).
\]
令每条边 \(e:i\to j\) 除携带 \(\gamma(e)\in\widehat{\ZZ}\) 外，再携带一个通道增量 \(\eta(e)\in G_{\mathrm{chan}}\)。
则对应的提升算子可写为
\[
(T_{\mathrm{GUT}}f)_j(x,h)
:=\sum_{e:i\to j} w(e)\,f_i\!\bigl(x-\gamma(e),\ h\cdot \eta(e)\bigr),
\qquad (x,h)\in \widehat{\ZZ}\times G_{\mathrm{chan}}.
\]
在对偶侧，\(\widehat{\ZZ}\) 的角色分解与 \(G_{\mathrm{chan}}\) 的 Peter--Weyl 分解把它同时块对角化为以 \((\chi,\varrho)\) 标记的有限维扭曲块：
\[
T_{\mathrm{GUT}}\ \simeq\ \bigoplus_{\chi\in\widehat{\ZZ}}\ \bigoplus_{\varrho\in\widehat{G_{\mathrm{chan}}}}\ M_{\chi,\varrho},
\qquad
(M_{\chi,\varrho})_{ij}:=\sum_{e:i\to j} w(e)\,\chi(\gamma(e))\,\varrho(\eta(e)).
\]
于是“大统一”可以在完全有限维、可枚举的谱数据层面表述为：对一组被观测门激活的通道族 \((\chi,\varrho)\)，其谱隙耦合 \(g_{\chi,\varrho}(u)\) 在某一统一点上发生对齐（见小节 \ref{subsec:group-unification-spectral-alignment} 的 minimax 判据）。
\end{remark}

\begin{remark}[Adelic 异常泛函：局部--整体异常抵消的可证伪统一判据]\label{rem:pom-adelic-anom-functional}
谱隙对齐（小节 \ref{subsec:group-unification-spectral-alignment}）给出的是指数级慢模的统一口径；
若希望把“\(E_m\)--\(\mathrm{LIFT}\)”不可交换的\emph{常数级残差}也纳入同一可计算框架，
可在定义 \ref{def:pom-anom-signature} 的异常签名 \(\mathsf{Anom}_G(K;\theta)\) 之上，对每个有限层（例如素数模轴 \(G_p=\ZZ/p\ZZ\)）引入一个标量化的局部异常强度。
一种对优化/梯度更友好的选择是
\[
A_p(\theta):=\bigl\|\mathsf{Anom}_{G_p}(K;\theta)\bigr\|_{\ell^2}^2
=\sum_{\chi\in\widehat{G_p}\setminus\{\chi_0\}}\bigl|\log\mathfrak{M}_\chi(\theta)\bigr|^2,
\]
其中 \(\chi_0\) 为平凡角色（亦可改用 \(\ell^\infty\) 范数 \(\max_{\chi\ne\chi_0}|\log\mathfrak{M}_\chi(\theta)|\) 作为审计口径）。
在 profinite 口径 \(G_\infty=\widehat{\ZZ}\) 下取素数截断 \(P\)，定义整体异常泛函
\[
\mathcal{A}_{\le P}(\theta):=\sum_{p\le P} w_p\,A_p(\theta),
\]
其中权重 \(w_p\) 可取 \(\log p\)、\(p^{-2}\) 等以反映尺度缩放或通道预算。
于是“跨模数的一致提升/异常抵消”可被表述为一个完全可证伪、且可由脚本直接审计的统一判据：
若存在 \(\theta^\star\) 使 \(\mathcal{A}_{\le P}(\theta)\) 在 \(\theta^\star\) 附近达到稳定极小（或满足 \(\nabla_\theta\mathcal{A}_{\le P}(\theta^\star)\approx 0\)），
并且随着截断 \(P\) 增大该最小点及最小值在数值上保持稳定（对截断深度/采样策略鲁棒），
则可把 \(\theta^\star\) 视为一个“局部--整体的 adelic 平坦点”（即常数级曲率/异常的整体抵消）。
由于每个 \(A_p(\theta)\) 都可由有限维扭曲块 \(M_\chi\) 的 \(\zeta\)/有限部接口计算得到，上述泛函与稳定性条件可被直接转化为可重复的优化与回归证书。
\end{remark}

\paragraph{归一化谱与 GRH 语义}
记平凡通道的 Perron 根为 \(\lambda:=\rho(M_{\mathbf{1}})\)，并对任意角色 \(\chi\) 定义归一化纤维算子
\[
U_\chi:=\lambda^{-1/2}M_\chi.
\]
令 \(w_{\chi,1},\dots,w_{\chi,|V|}\) 为 \(U_\chi\) 的特征值，则“临界圆”条件为 \(|w_{\chi,i}|=1\)。
由 \S\ref{subsec:zeta-online-kernel} 的谱--\(s\) 平面字典（命题 \ref{prop:s-plane-spectrum-poles}）以及附录 \ref{app:real-input-40-finite-rh} 的 toy-RH 讨论可知，这等价于把对应极点压到 \(\Re(s)=\tfrac12\)。
因此核版 GRH 的一个可操作“谱酉化指标”可取
\[
\mathcal{E}(\chi):=\max_{1\le i\le |V|}\bigl|\log|w_{\chi,i}|\bigr|,
\]
并把“GRH-型断言”理解为：在所关心的角色域/尺度域内，\(\mathcal{E}(\chi)\) 足够小（或在极限中趋于零）。

\paragraph{双尺度障碍与扩散口径}
当把角色族扩展到 \(\QQ/\ZZ\)（包含任意高阶小角角色）时，“最坏扭曲谱半径比”一般不可能在全部模数上保持统一远离 \(1\) 的界。
同步核加权情形里，附录已给出 near-coboundary 小角二次律并导出扩散阈值 \(n\asymp m^2\)（推论 \ref{cor:sync-rho-gap-m8-diffusive}）。
这一现象提示：在 profinite 逆极限语义下，更自然的 GRH/Chebotarev 形式应采用长度 \(n\) 与模数 \(m\) 的双极限尺度（参见猜想 \ref{conj:scaled-kernel-grh-diffusive} 的抽象表述）。
